% Comparación trasladada detrás de la ley y de la ontología sectorial.
% Se excluyen únicamente la portada autónoma y el metalenguaje de edición.
\section{Évaluation quantitative postérieure à la fixation du graphe génératif}
Le tableau suivant place dans des colonnes contiguës toute coordonnée scalaire que la théorie permet d’évaluer de manière reproductible. La référence expérimentale n’est consultée qu’après fixation de la sortie HMT--MD. Une ligne intitulée évaluation conditionnelle énonce exactement la valeur de la fonction une fois la signature fixée ; elle n’attribue pas à cette signature un sélecteur physique qui n’aurait pas été construit. Les quarks sont présentés sans résidu, car une comparaison numérique exige la concordance du schéma et de l’échelle de renormalisation.

Le dénombrement numérique de ce tableau est exhaustif relativement aux coordonnées scalaires actuellement reproductibles : un électron bêta fermé, dix-sept évaluations conditionnelles --- onze élémentaires et six composées --- et deux sorties nulles exactes. Il contient donc vingt lignes numériques. Ce nombre n’est pas le domaine de la loi. Les classes dont la sortie naturelle demeure un opérateur, un spectre, une orbite ou une multisection apparaissent ensuite avec cet objet complet, sans remplir leur colonne par la référence externe. L’inventaire expérimental de 471 masses et de 384 largeurs est conservé intégralement comme comparaison état par état ; son cardinal n’est pas présenté comme un dénombrement des prédictions HMT.

\Needspace{7\baselineskip}
\begingroup
\footnotesize
\setlength{\tabcolsep}{2.1pt}
\begin{longtable}{L{1.17cm}L{2.65cm}L{3.12cm}L{2.18cm}L{5.69cm}}
\toprule
\rowcolor{HMTNavy}\HMTtablehead{État} & \HMTtablehead{Valeur HMT reproductible (MeV/\allowbreak{}c\({}^{2}\))} & \HMTtablehead{Referencia posterior} & \HMTtablehead{Différence relative (ppm)} & \HMTtablehead{Nature et portée} \\
\midrule
\endfirsthead
\multicolumn{5}{l}{\sffamily\scriptsize\color{HMTGray}Suite du tableau}\\[-1pt]
\toprule
\rowcolor{HMTNavy}\HMTtablehead{État} & \HMTtablehead{Valeur HMT reproductible (MeV/\allowbreak{}c\({}^{2}\))} & \HMTtablehead{Referencia posterior} & \HMTtablehead{Différence relative (ppm)} & \HMTtablehead{Nature et portée} \\
\midrule
\endhead
\midrule
\endfoot
\bottomrule
\endlastfoot
e-/\allowbreak{}e+ & 0.510998\allowbreak{}950690\allowbreak{}000809 & 0.510998\allowbreak{}95069±\allowbreak{}0.000000\allowbreak{}00016 MeV/\allowbreak{}c\({}^{2}\) & 1.582406\allowbreak{}884×10\({}^{-9}\) & coordonnée bêta fermée dans la fibre centrale \\
\rowcolor{HMTLight}μ-/\allowbreak{}μ+ & 105.658360\allowbreak{}294435\allowbreak{}829303 & 105.658375\allowbreak{}5±\allowbreak{}0.000002\allowbreak{}3 MeV/\allowbreak{}c\({}^{2}\) & -0.143912\allowbreak{}53 & raffinement E3 conditionné par la signature \\
τ-/\allowbreak{}τ+ & 1776.860917\allowbreak{}631432\allowbreak{}295595 & 1776.932465\allowbreak{}134090\allowbreak{}1 MeV/\allowbreak{}c\({}^{2}\) (valeur centrale utilisée ; présentation arrondie : 1776.93±\allowbreak{}0.09 MeV/\allowbreak{}c\({}^{2}\)) & -40.264615\allowbreak{}601 & raffinement E3 conditionné par la signature \\
\rowcolor{HMTLight}u/\allowbreak{}\(\bar{\mathrm{u}}\) & 2.165924\allowbreak{}536173\allowbreak{}907663 & 2.16±\allowbreak{}0.07 MeV/\allowbreak{}c\({}^{2}\) & --- & carte angulaire ; aucune signature Z n’a été déclarée\({}^{5}\) ni aucun raffinement E3 \\
d/\allowbreak{}\(\bar{\mathrm{d}}\) & 4.679550\allowbreak{}162948\allowbreak{}874172 & 4.70±\allowbreak{}0.07 MeV/\allowbreak{}c\({}^{2}\) & --- & carte angulaire ; aucune signature Z n’a été déclarée\({}^{5}\) ni aucun raffinement E3 \\
\rowcolor{HMTLight}c/\allowbreak{}\(\bar{\mathrm{c}}\) & 1270.500838\allowbreak{}641911\allowbreak{}135286 & 1.2729±\allowbreak{}0.0045 GeV/\allowbreak{}c\({}^{2}\) & --- & carte angulaire ; aucune signature Z n’a été déclarée\({}^{5}\) ni aucun raffinement E3 \\
s/\allowbreak{}\(\bar{\mathrm{s}}\) & 92.940093\allowbreak{}907845\allowbreak{}640641 & 92.9 ±\allowbreak{}0.7 MeV/\allowbreak{}c\({}^{2}\) & --- & carte angulaire ; aucune signature Z n’a été déclarée\({}^{5}\) ni aucun raffinement E3 \\
\rowcolor{HMTLight}t/\allowbreak{}\(\bar{\mathrm{t}}\) & 172597.479544\allowbreak{}019136\allowbreak{}261367 & 172.60±\allowbreak{}0.27 GeV/\allowbreak{}c\({}^{2}\) & --- & carte angulaire ; aucune signature Z n’a été déclarée\({}^{5}\) ni aucun raffinement E3 \\
b/\allowbreak{}\(\bar{\mathrm{b}}\) & 4190.119763\allowbreak{}299790\allowbreak{}218075 & 4.186±\allowbreak{}0.006 GeV/\allowbreak{}c\({}^{2}\) & --- & carte angulaire ; aucune signature Z n’a été déclarée\({}^{5}\) ni aucun raffinement E3 \\
\rowcolor{HMTLight}γ & 0 & <1×10\({}^{-18}\) eV/\allowbreak{}c\({}^{2}\) & --- & sortie nulle exacte \\
g & 0 & --- & --- & sortie nulle exacte \\
\rowcolor{HMTLight}W+/\allowbreak{}W- & 80378.964909\allowbreak{}704547\allowbreak{}024865 & 80.3625 ±\allowbreak{}0.0077 GeV/\allowbreak{}c\({}^{2}\) & 204.882995\allowbreak{}235 & raffinement E3 conditionné par la signature \\
Z0 & 91187.533444\allowbreak{}313407\allowbreak{}844855 & 91.187873\allowbreak{}297221\allowbreak{}344 GeV/\allowbreak{}c\({}^{2}\) (valeur centrale utilisée ; présentation arrondie : 91.1879±\allowbreak{}0.0020 GeV/\allowbreak{}c\({}^{2}\)) & -3.726952\allowbreak{}89 & raffinement E3 conditionné par la signature \\
\rowcolor{HMTLight}H & 125292.466042\allowbreak{}536133\allowbreak{}000433 & 125.130943\allowbreak{}828161\allowbreak{}51 GeV/\allowbreak{}c\({}^{2}\) (valeur centrale utilisée ; présentation arrondie : 125.13±\allowbreak{}0.11 GeV/\allowbreak{}c\({}^{2}\)) & 1290.825509\allowbreak{}927 & carte angulaire ; aucune signature Z n’a été déclarée\({}^{5}\) ni aucun raffinement E3 \\
π\({}^{0}\) & 134.976724\allowbreak{}327872\allowbreak{}125739 & 134.976827\allowbreak{}767685 MeV/\allowbreak{}c\({}^{2}\) & -0.766352\allowbreak{}378 & raffinement E3 conditionné par la signature composée \\
\rowcolor{HMTLight}π±\allowbreak{} & 139.570485\allowbreak{}171572\allowbreak{}191375 & 139.570390\allowbreak{}983681 MeV/\allowbreak{}c\({}^{2}\) & 0.674841\allowbreak{}494 & raffinement E3 conditionné par la signature composée \\
K±\allowbreak{} & 493.677085\allowbreak{}649530\allowbreak{}612476 & 493.676599\allowbreak{}458041 MeV/\allowbreak{}c\({}^{2}\) & 0.984838\allowbreak{}03 & raffinement E3 conditionné par la signature composée \\
\rowcolor{HMTLight}K\({}^{0}\) & 497.611186\allowbreak{}497750\allowbreak{}457359 & 497.610767\allowbreak{}531258 MeV/\allowbreak{}c\({}^{2}\) & 0.841956\allowbreak{}243 & raffinement E3 conditionné par la signature composée \\
proton & 938.271751\allowbreak{}546984\allowbreak{}898299 & 938.272089\allowbreak{}43 MeV/\allowbreak{}c\({}^{2}\) & -0.360111\allowbreak{}975 & raffinement E3 conditionné par la signature composée \\
\rowcolor{HMTLight}neutron & 939.565810\allowbreak{}472507\allowbreak{}023313 & 939.565421\allowbreak{}94 MeV/\allowbreak{}c\({}^{2}\) & 0.413523\allowbreak{}633 & raffinement E3 conditionné par la signature composée \\
\end{longtable}
\endgroup
\Needspace{7\baselineskip}
\section{Évaluation individuelle de \(20\) sections matérielles et traçabilité de leurs coordonnées métrologiques}
\Needspace{7\baselineskip}
\subsection{Paire électron–positron : parcours central, conjugaison des feuillets et caractère de masse}
\begin{itemize}
\item \textbf{Identifiant du catalogue :} SM-LC-1.
\item \textbf{Type de descripteur :} sélecteur basal et lecteurs bidirectionnels.
\item \textbf{Descripteur :} sélecteur (-22,+15); K\_D=\allowbreak{}K\_Ω=\allowbreak{}(80,54,6).
\item \textbf{Étape :} bêta bidirectionnel ; coordonnée initiale 0.510998\allowbreak{}922488\allowbreak{}066073 MeV/c\({}^{2}\).
\item \textbf{Évaluation HMT--MD reproductible :} 0.510998\allowbreak{}950690\allowbreak{}000809 MeV/c\({}^{2}\).
\item \textbf{Nature mathématique :} coordonnée bêta fermée dans la fibre centrale.
\item \textbf{Référence ultérieure :} 0.510998\allowbreak{}95069±\allowbreak{}0.000000\allowbreak{}00016 MeV/c\({}^{2}\).
\item \textbf{Identité et métrologie externes :} Particle Data Group (PDG); PDG2026:\allowbreak{}observable:\allowbreak{}S003:\allowbreak{}S003M:\allowbreak{}332838; schéma non conservé dans l’extrait documentaire; échelle non conservée dans l’extrait documentaire; niveau de confiance non conservé dans l’extrait documentaire.
\item \textbf{Différence relative:} 1.582406\allowbreak{}884×10\({}^{-9}\) ppm.
\item \textbf{Portée :} dérivation scalaire fermée.
\end{itemize}
\Needspace{7\baselineskip}
\subsection{Paire muon–antimuon : conjugaison des parcours et caractère leptonique}
\begin{itemize}
\item \textbf{Identifiant du catalogue :} SM-LC-2.
\item \textbf{Type de descripteur :} signature complète σ\(\in\)Z\({}^{5}\).
\item \textbf{Descripteur :} (42,-3,8,0,2).
\item \textbf{Étape :} tour bêta E3 ; coordonnée initiale 105.564059\allowbreak{}012859\allowbreak{}26472 MeV/c\({}^{2}\).
\item \textbf{Évaluation HMT--MD reproductible :} 105.658360\allowbreak{}294435\allowbreak{}829303 MeV/c\({}^{2}\).
\item \textbf{Nature mathématique :} raffinement E3 conditionné par la signature.
\item \textbf{Référence ultérieure :} 105.658375\allowbreak{}5±\allowbreak{}0.000002\allowbreak{}3 MeV/c\({}^{2}\).
\item \textbf{Identité et métrologie externes :} Particle Data Group (PDG); PDG2026:\allowbreak{}observable:\allowbreak{}S004:\allowbreak{}S004M:\allowbreak{}332849; schéma non conservé dans l’extrait documentaire; échelle non conservée dans l’extrait documentaire; niveau de confiance non conservé dans l’extrait documentaire.
\item \textbf{Différence relative:} -0.143912\allowbreak{}53 ppm.
\item \textbf{Portée :} évaluation exacte une fois la signature fixée ; la sélection physique de cette signature constitue une opération indépendante.
\end{itemize}
\Needspace{7\baselineskip}
\subsection{Paire tau–antitau : conjugaison des parcours et caractère leptonique}
\begin{itemize}
\item \textbf{Identifiant du catalogue :} SM-LC-3.
\item \textbf{Type de descripteur :} signature complète σ\(\in\)Z\({}^{5}\).
\item \textbf{Descripteur :} (64,0,25,-1,6).
\item \textbf{Étape :} tour bêta E3 ; coordonnée initiale 1771.945393\allowbreak{}886234\allowbreak{}110148 MeV/c\({}^{2}\).
\item \textbf{Évaluation HMT--MD reproductible :} 1776.860917\allowbreak{}631432\allowbreak{}295595 MeV/c\({}^{2}\).
\item \textbf{Nature mathématique :} raffinement E3 conditionné par la signature.
\item \textbf{Référence ultérieure :} 1776.932465\allowbreak{}134090\allowbreak{}1 MeV/c\({}^{2}\) (valeur centrale utilisée ; présentation arrondie : 1776.93±\allowbreak{}0.09 MeV/c\({}^{2}\)).
\item \textbf{Identité et métrologie externes :} Particle Data Group (PDG); PDG2026:\allowbreak{}observable:\allowbreak{}S035:\allowbreak{}S035M:\allowbreak{}332893; schéma non conservé dans l’extrait documentaire; échelle non conservée dans l’extrait documentaire; niveau de confiance non conservé dans l’extrait documentaire.
\item \textbf{Différence relative:} -40.264615\allowbreak{}601 ppm.
\item \textbf{Portée :} évaluation exacte une fois la signature fixée ; la sélection physique de cette signature constitue une opération indépendante.
\end{itemize}
\Needspace{7\baselineskip}
\subsection{Quark haut et antiquark haut}
\begin{itemize}
\item \textbf{Identifiant du catalogue :} SM-Q-U1.
\item \textbf{Type de descripteur :} coindexation angulaire (n\_A,s\_C)\(\in\)Z\({}^{2}\).
\item \textbf{Descripteur :} (11,7).
\item \textbf{Étape :} carte angulaire conditionnelle ; coordonnée initiale 2.165924\allowbreak{}536173\allowbreak{}907663 MeV/c\({}^{2}\).
\item \textbf{Évaluation HMT--MD reproductible :} 2.165924\allowbreak{}536173\allowbreak{}907663 MeV/c\({}^{2}\).
\item \textbf{Nature mathématique :} carte angulaire ; aucune signature Z n’a été déclarée\({}^{5}\) ni aucun raffinement E3.
\item \textbf{Référence ultérieure :} 2.16±\allowbreak{}0.07 MeV/c\({}^{2}\).
\item \textbf{Identité et métrologie externes :} Particle Data Group (PDG); PDG2026:\allowbreak{}observable:\allowbreak{}Q002:\allowbreak{}Q002M:\allowbreak{}333588; schéma non conservé dans l’extrait documentaire; échelle non conservée dans l’extrait documentaire; niveau de confiance 90\%.
\item \textbf{Différence relative:} --- ppm.
\item \textbf{Portée :} évaluation conditionnelle ; la signature n’a pas été choisie à partir de la valeur comparée ; schéma et échelle externes requis pour le résidu.
\end{itemize}
\Needspace{7\baselineskip}
\subsection{Quark bas et antiquark bas}
\begin{itemize}
\item \textbf{Identifiant du catalogue :} SM-Q-D1.
\item \textbf{Type de descripteur :} coindexation angulaire (n\_A,s\_C)\(\in\)Z\({}^{2}\).
\item \textbf{Descripteur :} (17,8).
\item \textbf{Étape :} carte angulaire conditionnelle ; coordonnée initiale 4.679550\allowbreak{}162948\allowbreak{}874172 MeV/c\({}^{2}\).
\item \textbf{Évaluation HMT--MD reproductible :} 4.679550\allowbreak{}162948\allowbreak{}874172 MeV/c\({}^{2}\).
\item \textbf{Nature mathématique :} carte angulaire ; aucune signature Z n’a été déclarée\({}^{5}\) ni aucun raffinement E3.
\item \textbf{Référence ultérieure :} 4.70±\allowbreak{}0.07 MeV/c\({}^{2}\).
\item \textbf{Identité et métrologie externes :} Particle Data Group (PDG); PDG2026:\allowbreak{}observable:\allowbreak{}Q001:\allowbreak{}Q001M:\allowbreak{}333589; schéma non conservé dans l’extrait documentaire; échelle non conservée dans l’extrait documentaire; niveau de confiance 90\%.
\item \textbf{Différence relative:} --- ppm.
\item \textbf{Portée :} évaluation conditionnelle ; la signature n’a pas été choisie à partir de la valeur comparée ; schéma et échelle externes requis pour le résidu.
\end{itemize}
\Needspace{7\baselineskip}
\subsection{Quark charme et antiquark charme}
\begin{itemize}
\item \textbf{Identifiant du catalogue :} SM-Q-U2.
\item \textbf{Type de descripteur :} coindexation angulaire (n\_A,s\_C)\(\in\)Z\({}^{2}\).
\item \textbf{Descripteur :} (61,8).
\item \textbf{Étape :} carte angulaire conditionnelle ; coordonnée initiale 1270.500838\allowbreak{}641911\allowbreak{}135286 MeV/c\({}^{2}\).
\item \textbf{Évaluation HMT--MD reproductible :} 1270.500838\allowbreak{}641911\allowbreak{}135286 MeV/c\({}^{2}\).
\item \textbf{Nature mathématique :} carte angulaire ; aucune signature Z n’a été déclarée\({}^{5}\) ni aucun raffinement E3.
\item \textbf{Référence ultérieure :} 1.2729±\allowbreak{}0.0045 GeV/c\({}^{2}\).
\item \textbf{Identité et métrologie externes :} Particle Data Group (PDG); PDG2026:\allowbreak{}observable:\allowbreak{}Q004:\allowbreak{}Q004M:\allowbreak{}333604; schéma non conservé dans l’extrait documentaire; échelle non conservée dans l’extrait documentaire; niveau de confiance 90\%.
\item \textbf{Différence relative:} --- ppm.
\item \textbf{Portée :} évaluation conditionnelle ; la signature n’a pas été choisie à partir de la valeur comparée ; schéma et échelle externes requis pour le résidu.
\end{itemize}
\Needspace{7\baselineskip}
\subsection{Quark étrange et antiquark étrange}
\begin{itemize}
\item \textbf{Identifiant du catalogue :} SM-Q-D2.
\item \textbf{Type de descripteur :} coindexation angulaire (n\_A,s\_C)\(\in\)Z\({}^{2}\).
\item \textbf{Descripteur :} (41,-3).
\item \textbf{Étape :} carte angulaire conditionnelle ; coordonnée initiale 92.940093\allowbreak{}907845\allowbreak{}640641 MeV/c\({}^{2}\).
\item \textbf{Évaluation HMT--MD reproductible :} 92.940093\allowbreak{}907845\allowbreak{}640641 MeV/c\({}^{2}\).
\item \textbf{Nature mathématique :} carte angulaire ; aucune signature Z n’a été déclarée\({}^{5}\) ni aucun raffinement E3.
\item \textbf{Référence ultérieure :} 92.9 ±\allowbreak{}0.7 MeV/c\({}^{2}\).
\item \textbf{Identité et métrologie externes :} Particle Data Group (PDG); PDG2026:\allowbreak{}observable:\allowbreak{}Q003:\allowbreak{}Q003M:\allowbreak{}333590; schéma non conservé dans l’extrait documentaire; échelle non conservée dans l’extrait documentaire; niveau de confiance 90\%.
\item \textbf{Différence relative:} --- ppm.
\item \textbf{Portée :} évaluation conditionnelle ; la signature n’a pas été choisie à partir de la valeur comparée ; schéma et échelle externes requis pour le résidu.
\end{itemize}
\Needspace{7\baselineskip}
\subsection{Quark sommet et antiquark sommet}
\begin{itemize}
\item \textbf{Identifiant du catalogue :} SM-Q-U3.
\item \textbf{Type de descripteur :} coindexation angulaire (n\_A,s\_C)\(\in\)Z\({}^{2}\).
\item \textbf{Descripteur :} (100,-1).
\item \textbf{Étape :} carte angulaire conditionnelle ; coordonnée initiale 172597.479544\allowbreak{}019136\allowbreak{}261367 MeV/c\({}^{2}\).
\item \textbf{Évaluation HMT--MD reproductible :} 172597.479544\allowbreak{}019136\allowbreak{}261367 MeV/c\({}^{2}\).
\item \textbf{Nature mathématique :} carte angulaire ; aucune signature Z n’a été déclarée\({}^{5}\) ni aucun raffinement E3.
\item \textbf{Référence ultérieure :} 172.60±\allowbreak{}0.27 GeV/c\({}^{2}\).
\item \textbf{Identité et métrologie externes :} Particle Data Group (PDG); PDG2026:\allowbreak{}observable:\allowbreak{}Q007:\allowbreak{}Q007TP:\allowbreak{}333614; schéma non conservé dans l’extrait documentaire; échelle non conservée dans l’extrait documentaire; niveau de confiance non conservé dans l’extrait documentaire.
\item \textbf{Différence relative:} --- ppm.
\item \textbf{Portée :} évaluation conditionnelle ; la signature n’a pas été choisie à partir de la valeur comparée ; schéma et échelle externes requis pour le résidu.
\end{itemize}
\Needspace{7\baselineskip}
\subsection{Quark fond et antiquark fond}
\begin{itemize}
\item \textbf{Identifiant du catalogue :} SM-Q-D3.
\item \textbf{Type de descripteur :} coindexation angulaire (n\_A,s\_C)\(\in\)Z\({}^{2}\).
\item \textbf{Descripteur :} (71,-5).
\item \textbf{Étape :} carte angulaire conditionnelle ; coordonnée initiale 4190.119763\allowbreak{}299790\allowbreak{}218075 MeV/c\({}^{2}\).
\item \textbf{Évaluation HMT--MD reproductible :} 4190.119763\allowbreak{}299790\allowbreak{}218075 MeV/c\({}^{2}\).
\item \textbf{Nature mathématique :} carte angulaire ; aucune signature Z n’a été déclarée\({}^{5}\) ni aucun raffinement E3.
\item \textbf{Référence ultérieure :} 4.186±\allowbreak{}0.006 GeV/c\({}^{2}\).
\item \textbf{Identité et métrologie externes :} Particle Data Group (PDG); PDG2026:\allowbreak{}observable:\allowbreak{}Q005:\allowbreak{}Q005M:\allowbreak{}333611; schéma non conservé dans l’extrait documentaire; échelle non conservée dans l’extrait documentaire; niveau de confiance 90\%.
\item \textbf{Différence relative:} --- ppm.
\item \textbf{Portée :} évaluation conditionnelle ; la signature n’a pas été choisie à partir de la valeur comparée ; schéma et échelle externes requis pour le résidu.
\end{itemize}
\Needspace{7\baselineskip}
\subsection{Descripteur spectral reproductible du photon : secteur nul, hélicité et caractère électromagnétique}
\begin{itemize}
\item \textbf{Identifiant du catalogue :} SM-GA-EM.
\item \textbf{Type de descripteur :} invariant sectoriel nul.
\item \textbf{Descripteur :} carte nulle.
\item \textbf{Étape :} carte nulle ; coordonnée initiale 0 MeV/c\({}^{2}\).
\item \textbf{Évaluation HMT--MD reproductible :} 0 MeV/c\({}^{2}\).
\item \textbf{Nature mathématique :} sortie nulle exacte.
\item \textbf{Référence ultérieure :} <1×10\({}^{-18}\) eV/c\({}^{2}\).
\item \textbf{Identité et métrologie externes :} Particle Data Group (PDG); PDG2026:\allowbreak{}observable:\allowbreak{}S000:\allowbreak{}S000M:\allowbreak{}332461; schéma non conservé dans l’extrait documentaire; échelle non conservée dans l’extrait documentaire; niveau de confiance non conservé dans l’extrait documentaire.
\item \textbf{Différence relative:} --- ppm.
\item \textbf{Portée :} exact dans le secteur nul ; la borne externe n’est pas interprétée comme une masse positive.
\item \textbf{Borne photonique :} référence externe illustrative non comparable ; l’extrait identifie l’observable, mais ne conserve pas le niveau de confiance.
\end{itemize}
\Needspace{7\baselineskip}
\subsection{Réalisation du gluon comme section de la représentation adjointe de \(SU(3)\)}
\begin{itemize}
\item \textbf{Identifiant du catalogue :} SM-GA-GL.
\item \textbf{Type de descripteur :} invariant sectoriel nul.
\item \textbf{Descripteur :} carte nulle.
\item \textbf{Étape :} carte nulle ; coordonnée initiale 0 MeV/c\({}^{2}\).
\item \textbf{Évaluation HMT--MD reproductible :} 0 MeV/c\({}^{2}\).
\item \textbf{Nature mathématique :} sortie nulle exacte.
\item \textbf{Referencia posterior:} ---.
\item \textbf{Identité et métrologie externes :} source non consignée ; observable non individualisé ; schéma non consigné ; échelle non consignée ; niveau de confiance non consigné.
\item \textbf{Différence relative:} --- ppm.
\item \textbf{Portée :} exact dans le secteur nul ; la nullité appartient à la carte de jauge de couleur.
\end{itemize}
\Needspace{7\baselineskip}
\subsection{Bosons \texorpdfstring{\(W^{+}\)}{W+} et \texorpdfstring{\(W^{-}\)}{W−}}
\begin{itemize}
\item \textbf{Identifiant du catalogue :} SM-GA-W.
\item \textbf{Type de descripteur :} signature complète σ\(\in\)Z\({}^{5}\).
\item \textbf{Descripteur :} (94,-1,0,-2,4).
\item \textbf{Étape :} tour bêta E3 ; coordonnée initiale 80381.592550\allowbreak{}220666\allowbreak{}304144 MeV/c\({}^{2}\).
\item \textbf{Évaluation HMT--MD reproductible :} 80378.964909\allowbreak{}704547\allowbreak{}024865 MeV/c\({}^{2}\).
\item \textbf{Nature mathématique :} raffinement E3 conditionné par la signature.
\item \textbf{Référence ultérieure :} 80.3625 ±\allowbreak{}0.0077 GeV/c\({}^{2}\).
\item \textbf{Identité et métrologie externes :} Particle Data Group (PDG); PDG2026:\allowbreak{}observable:\allowbreak{}S043:\allowbreak{}S043M:\allowbreak{}332464; schéma non conservé dans l’extrait documentaire; échelle non conservée dans l’extrait documentaire; niveau de confiance non conservé dans l’extrait documentaire.
\item \textbf{Différence relative:} 204.882995\allowbreak{}235 ppm.
\item \textbf{Portée :} évaluation exacte une fois la signature fixée ; la sélection physique de cette signature constitue une opération indépendante.
\item \textbf{Convention résonante :} l’extrait externe ne fixe pas conjointement la convention de pôle, de Breit--Wigner ou sur couche de masse, la largeur et la covariance ; le résidu n’est interprété que comme information scalaire.
\end{itemize}
\Needspace{7\baselineskip}
\subsection{Descripteur spectral reproductible du boson \(Z^0\) : mélange neutre, polarisation et caractère électrofaible}
\begin{itemize}
\item \textbf{Identifiant du catalogue :} SM-GA-Z.
\item \textbf{Type de descripteur :} signature complète σ\(\in\)Z\({}^{5}\).
\item \textbf{Descripteur :} (95,-1,-11,0,-4).
\item \textbf{Étape :} tour bêta E3 ; coordonnée initiale 91299.748286\allowbreak{}598170\allowbreak{}117317 MeV/c\({}^{2}\).
\item \textbf{Évaluation HMT--MD reproductible :} 91187.533444\allowbreak{}313407\allowbreak{}844855 MeV/c\({}^{2}\).
\item \textbf{Nature mathématique :} raffinement E3 conditionné par la signature.
\item \textbf{Référence ultérieure :} 91.187873\allowbreak{}297221\allowbreak{}344 GeV/c\({}^{2}\) (valeur centrale utilisée ; présentation arrondie : 91.1879±\allowbreak{}0.0020 GeV/c\({}^{2}\)).
\item \textbf{Identité et métrologie externes :} Particle Data Group (PDG); PDG2026:\allowbreak{}observable:\allowbreak{}S044:\allowbreak{}S044M:\allowbreak{}332513; schéma non conservé dans l’extrait documentaire; échelle non conservée dans l’extrait documentaire; niveau de confiance non conservé dans l’extrait documentaire.
\item \textbf{Différence relative:} -3.726952\allowbreak{}89 ppm.
\item \textbf{Portée :} évaluation exacte une fois la signature fixée ; la sélection physique de cette signature constitue une opération indépendante.
\item \textbf{Convention résonante :} l’extrait externe ne fixe pas conjointement la convention de pôle, de Breit--Wigner ou sur couche de masse, la largeur et la covariance ; le résidu n’est interprété que comme information scalaire.
\end{itemize}
\Needspace{7\baselineskip}
\subsection{Réalisation du boson de Higgs : mode scalaire, couplage et caractère de masse}
\begin{itemize}
\item \textbf{Identifiant du catalogue :} SM-H-1.
\item \textbf{Type de descripteur :} coindexation angulaire (n\_A,s\_C)\(\in\)Z\({}^{2}\).
\item \textbf{Descripteur :} (97,9).
\item \textbf{Étape :} carte angulaire conditionnelle ; coordonnée initiale 125292.466042\allowbreak{}536133\allowbreak{}000433 MeV/c\({}^{2}\).
\item \textbf{Évaluation HMT--MD reproductible :} 125292.466042\allowbreak{}536133\allowbreak{}000433 MeV/c\({}^{2}\).
\item \textbf{Nature mathématique :} carte angulaire ; aucune signature Z n’a été déclarée\({}^{5}\) ni aucun raffinement E3.
\item \textbf{Référence ultérieure :} 125.130943\allowbreak{}828161\allowbreak{}51 GeV/c\({}^{2}\) (valeur centrale utilisée ; présentation arrondie : 125.13±\allowbreak{}0.11 GeV/c\({}^{2}\)).
\item \textbf{Identité et métrologie externes :} Particle Data Group (PDG); PDG2026:\allowbreak{}observable:\allowbreak{}S126:\allowbreak{}S126M:\allowbreak{}332733; schéma non conservé dans l’extrait documentaire; échelle non conservée dans l’extrait documentaire; niveau de confiance non conservé dans l’extrait documentaire.
\item \textbf{Différence relative:} 1290.825509\allowbreak{}927 ppm.
\item \textbf{Portée :} évaluation conditionnelle ; la signature n’a pas été choisie à partir de la valeur comparée.
\item \textbf{Convention résonante :} l’extrait externe ne fixe pas conjointement la convention de pôle, de Breit--Wigner ou sur couche de masse, la largeur et la covariance ; le résidu n’est interprété que comme information scalaire.
\end{itemize}
\Needspace{7\baselineskip}
\subsection{Pion neutre \texorpdfstring{\(\pi^{0}\)}{pi0}}
\begin{itemize}
\item \textbf{Identifiant du catalogue :} CMP-PI0.
\item \textbf{Type de descripteur :} signature complète σ\(\in\)Z\({}^{5}\).
\item \textbf{Descripteur :} (44,-4,-25,1,-2).
\item \textbf{Étape :} tour bêta E3 composée ; coordonnée initiale 135.350257\allowbreak{}251504\allowbreak{}283349 MeV/c\({}^{2}\).
\item \textbf{Évaluation HMT--MD reproductible :} 134.976724\allowbreak{}327872\allowbreak{}125739 MeV/c\({}^{2}\).
\item \textbf{Nature mathématique :} raffinement E3 conditionné par la signature composée.
\item \textbf{Référence ultérieure :} 134.976827\allowbreak{}767685 MeV/c\({}^{2}\).
\item \textbf{Identité et métrologie externes :} référence composée conservée par le supplément E3 ; la convention, la covariance et l’édition sont consignées dans la source métrologique de référence.
\item \textbf{Différence relative:} -0.766352\allowbreak{}378 ppm.
\item \textbf{Portée :} évaluation exacte une fois la signature composée fixée ; le générateur de liaison demeure un morphisme propre.
\end{itemize}
\Needspace{7\baselineskip}
\subsection{Pions chargés \(\pi^\pm\) : composition quark–antiquark, charge et masse}
\begin{itemize}
\item \textbf{Identifiant du catalogue :} CMP-PIPM.
\item \textbf{Type de descripteur :} signature complète σ\(\in\)Z\({}^{5}\).
\item \textbf{Descripteur :} (44,1,-2,2,-1).
\item \textbf{Étape :} tour bêta E3 composée ; coordonnée initiale 139.596265\allowbreak{}529971\allowbreak{}781015 MeV/c\({}^{2}\).
\item \textbf{Évaluation HMT--MD reproductible :} 139.570485\allowbreak{}171572\allowbreak{}191375 MeV/c\({}^{2}\).
\item \textbf{Nature mathématique :} raffinement E3 conditionné par la signature composée.
\item \textbf{Référence ultérieure :} 139.570390\allowbreak{}983681 MeV/c\({}^{2}\).
\item \textbf{Identité et métrologie externes :} référence composée conservée par le supplément E3 ; la convention, la covariance et l’édition sont consignées dans la source métrologique de référence.
\item \textbf{Différence relative:} 0.674841\allowbreak{}494 ppm.
\item \textbf{Portée :} évaluation exacte une fois la signature composée fixée ; le générateur de liaison demeure un morphisme propre.
\end{itemize}
\Needspace{7\baselineskip}
\subsection{Kaons chargés \(K^\pm\) : composition avec étrangeté, charge et masse}
\begin{itemize}
\item \textbf{Identifiant du catalogue :} CMP-KPM.
\item \textbf{Type de descripteur :} signature complète σ\(\in\)Z\({}^{5}\).
\item \textbf{Descripteur :} (54,-1,17,-2,2).
\item \textbf{Étape :} tour bêta E3 composée ; coordonnée initiale 492.762503\allowbreak{}210140\allowbreak{}547854 MeV/c\({}^{2}\).
\item \textbf{Évaluation HMT--MD reproductible :} 493.677085\allowbreak{}649530\allowbreak{}612476 MeV/c\({}^{2}\).
\item \textbf{Nature mathématique :} raffinement E3 conditionné par la signature composée.
\item \textbf{Référence ultérieure :} 493.676599\allowbreak{}458041 MeV/c\({}^{2}\).
\item \textbf{Identité et métrologie externes :} référence composée conservée par le supplément E3 ; la convention, la covariance et l’édition sont consignées dans la source métrologique de référence.
\item \textbf{Différence relative:} 0.984838\allowbreak{}03 ppm.
\item \textbf{Portée :} évaluation exacte une fois la signature composée fixée ; le générateur de liaison demeure un morphisme propre.
\end{itemize}
\Needspace{7\baselineskip}
\subsection{Kaon neutre \texorpdfstring{\(K^{0}\)}{K0}}
\begin{itemize}
\item \textbf{Identifiant du catalogue :} CMP-K0.
\item \textbf{Type de descripteur :} signature complète σ\(\in\)Z\({}^{5}\).
\item \textbf{Descripteur :} (54,0,32,3,-2).
\item \textbf{Étape :} tour bêta E3 composée ; coordonnée initiale 495.816066\allowbreak{}594426\allowbreak{}768595 MeV/c\({}^{2}\).
\item \textbf{Évaluation HMT--MD reproductible :} 497.611186\allowbreak{}497750\allowbreak{}457359 MeV/c\({}^{2}\).
\item \textbf{Nature mathématique :} raffinement E3 conditionné par la signature composée.
\item \textbf{Référence ultérieure :} 497.610767\allowbreak{}531258 MeV/c\({}^{2}\).
\item \textbf{Identité et métrologie externes :} référence composée conservée par le supplément E3 ; la convention, la covariance et l’édition sont consignées dans la source métrologique de référence.
\item \textbf{Différence relative:} 0.841956\allowbreak{}243 ppm.
\item \textbf{Portée :} évaluation exacte une fois la signature composée fixée ; le générateur de liaison demeure un morphisme propre.
\end{itemize}
\Needspace{7\baselineskip}
\subsection{Proton : composition baryonique, neutralisation de couleur et caractère de masse}
\begin{itemize}
\item \textbf{Identifiant du catalogue :} CMP-P.
\item \textbf{Type de descripteur :} signature complète σ\(\in\)Z\({}^{5}\).
\item \textbf{Descripteur :} (59,0,9,1,0).
\item \textbf{Étape :} tour bêta E3 composée ; coordonnée initiale 937.314779\allowbreak{}258699\allowbreak{}634563 MeV/c\({}^{2}\).
\item \textbf{Évaluation HMT--MD reproductible :} 938.271751\allowbreak{}546984\allowbreak{}898299 MeV/c\({}^{2}\).
\item \textbf{Nature mathématique :} raffinement E3 conditionné par la signature composée.
\item \textbf{Référence ultérieure :} 938.272089\allowbreak{}43 MeV/c\({}^{2}\).
\item \textbf{Identité et métrologie externes :} référence composée conservée par le supplément E3 ; la convention, la covariance et l’édition sont consignées dans la source métrologique de référence.
\item \textbf{Différence relative:} -0.360111\allowbreak{}975 ppm.
\item \textbf{Portée :} évaluation exacte une fois la signature composée fixée ; le générateur de liaison demeure un morphisme propre.
\end{itemize}
\Needspace{7\baselineskip}
\subsection{Neutron : composition baryonique, neutralisation de couleur et caractère de masse}
\begin{itemize}
\item \textbf{Identifiant du catalogue :} CMP-N.
\item \textbf{Type de descripteur :} signature complète σ\(\in\)Z\({}^{5}\).
\item \textbf{Descripteur :} (59,1,-34,0,1).
\item \textbf{Étape :} tour bêta E3 composée ; coordonnée initiale 943.123155\allowbreak{}648642\allowbreak{}014982 MeV/c\({}^{2}\).
\item \textbf{Évaluation HMT--MD reproductible :} 939.565810\allowbreak{}472507\allowbreak{}023313 MeV/c\({}^{2}\).
\item \textbf{Nature mathématique :} raffinement E3 conditionné par la signature composée.
\item \textbf{Référence ultérieure :} 939.565421\allowbreak{}94 MeV/c\({}^{2}\).
\item \textbf{Identité et métrologie externes :} référence composée conservée par le supplément E3 ; la convention, la covariance et l’édition sont consignées dans la source métrologique de référence.
\item \textbf{Différence relative:} 0.413523\allowbreak{}633 ppm.
\item \textbf{Portée :} évaluation exacte une fois la signature composée fixée ; le générateur de liaison demeure un morphisme propre.
\end{itemize}
% Las tablas históricas abreviadas permanecen conservadas a continuación en
% este propietario, pero quedan fuera del tronco científico vigente: no
% definen el dominio 81/56/13/324, la firma, el carácter ni las torres.
\endinput
\clearpage
\begin{landscape}
\section{Maillage historique étendu conservé}
\begingroup
\fontsize{6.8}{7.7}\selectfont
\setlength{\tabcolsep}{1.15pt}
\renewcommand{\arraystretch}{1.06}
\begin{longtable}{L{2.93cm}L{2.93cm}L{2.93cm}L{2.93cm}L{2.93cm}L{2.93cm}L{2.93cm}L{2.93cm}}
\toprule
\rowcolor{HMTNavy}\HMTtablehead{État} & \HMTtablehead{n} & \HMTtablehead{k} & \HMTtablehead{t} & \HMTtablehead{Évaluation HMT historique (MeV/\allowbreak{}c\({}^{2}\))} & \HMTtablehead{Référence ultérieure (MeV/\allowbreak{}c\({}^{2}\))} & \HMTtablehead{Différence consignée (ppm)} & \HMTtablehead{Force mathématique} \\
\midrule
\endfirsthead
\multicolumn{8}{l}{\sffamily\scriptsize\color{HMTGray}Suite du tableau}\\[-1pt]
\toprule
\rowcolor{HMTNavy}\HMTtablehead{État} & \HMTtablehead{n} & \HMTtablehead{k} & \HMTtablehead{t} & \HMTtablehead{Évaluation HMT historique (MeV/\allowbreak{}c\({}^{2}\))} & \HMTtablehead{Référence ultérieure (MeV/\allowbreak{}c\({}^{2}\))} & \HMTtablehead{Différence consignée (ppm)} & \HMTtablehead{Force mathématique} \\
\midrule
\endhead
\midrule
\endfoot
\bottomrule
\endlastfoot
mu & 64 & 146 & 1 & 105.656529 & 105.658374 & -17.461938\allowbreak{}2274 & preuve numérique historique conditionnelle ; ne constitue pas une règle physique de sélection antérieure à la comparaison \\
\rowcolor{HMTLight}tau & 64 & 349 & -1 & 1776.835 & 1776.86 & -14.069763\allowbreak{}5154 & preuve numérique historique conditionnelle ; ne constitue pas une règle physique de sélection antérieure à la comparaison \\
π±\allowbreak{} & 37 & 125 & 0 & 139.570582 & 139.57039 & 1.375649\allowbreak{}94982 & preuve numérique historique conditionnelle ; ne constitue pas une règle physique de sélection antérieure à la comparaison \\
\rowcolor{HMTLight}K±\allowbreak{} & 52 & 177 & 1 & 493.679 & 493.677 & 4.051231\allowbreak{}87833 & preuve numérique historique conditionnelle ; ne constitue pas une règle physique de sélection antérieure à la comparaison \\
K\({}^{0}\) & 53 & 178 & 1 & 497.616 & 497.611 & 10.048009\allowbreak{}3889 & preuve numérique historique conditionnelle ; ne constitue pas une règle physique de sélection antérieure à la comparaison \\
\rowcolor{HMTLight}p & 55 & 195 & 1 & 938.27244 & 938.272081 & 0.382618\allowbreak{}226919 & preuve numérique historique conditionnelle ; ne constitue pas une règle physique de sélection antérieure à la comparaison \\
n & 55 & 197 & 0 & 939.538 & 939.565413 & -29.176254\allowbreak{}9161 & preuve numérique historique conditionnelle ; ne constitue pas une règle physique de sélection antérieure à la comparaison \\
\rowcolor{HMTLight}eta & 46 & 158 & 0 & 547.836 & 547.862 & -47.457206\allowbreak{}3768 & preuve numérique historique conditionnelle ; ne constitue pas une règle physique de sélection antérieure à la comparaison \\
ρ\({}^{0}\) & 50 & 176 & 0 & 775.267 & 775.26 & 9.029228\allowbreak{}90385 & preuve numérique historique conditionnelle ; ne constitue pas une règle physique de sélection antérieure à la comparaison \\
\rowcolor{HMTLight}φ(1020) & 54 & 190 & 0 & 1019.475 & 1019.461 & 13.732747\allowbreak{}0104 & preuve numérique historique conditionnelle ; ne constitue pas une règle physique de sélection antérieure à la comparaison \\
J/\allowbreak{}ψ & 64 & 237 & 0 & 3096.979 & 3096.916 & 20.342818\allowbreak{}4684 & preuve numérique historique conditionnelle ; ne constitue pas une règle physique de sélection antérieure à la comparaison \\
\rowcolor{HMTLight}Υ(1S) & 89 & 334 & 0 & 9460.34 & 9460.3 & 4.228195\allowbreak{}72318 & preuve numérique historique conditionnelle ; ne constitue pas une règle physique de sélection antérieure à la comparaison \\
\end{longtable}
\endgroup
\end{landscape}
Ces douze valeurs ne sont pas fusionnées avec les vingt sorties du tableau de référence. Le maillage conserve cinq étiquettes supplémentaires ---η, ρ\({}^{0}\), φ(1020), J/ψ et Υ(1S)---, mais ne construit pas encore l’application état composé→parcours et enregistrement généalogique→(n,k,t). Le quantum interne CT108, 0.058177\allowbreak{}641733\allowbreak{}144319\allowbreak{}230789\allowbreak{}692282\allowbreak{}953757 dans la normalisation hbar\_int=\allowbreak{}c\_int=\allowbreak{}t0=\allowbreak{}1, est une section de la droite de masse et non une espèce.

Le tableau ne délimite pas le domaine. Le domaine interne demeure l’atlas de 81 cellules, 56 familles, 13 multisections et 324 parcours orientés ; les réalisations physiques typées et l’inventaire de 471 masses et de 384 largeurs sont développés ensuite.

\Needspace{7\baselineskip}
